\documentclass{article}
\usepackage{graphicx} % Required for inserting images

\title{30350 Remote Sensing Group Project - Regional noise reduction in NIMBUS 5 ESMR microwave radiometer sea ice concentration data}
\author{Mingming Li s252735 \\ Johannes Ruckdeschel s263265 \\ Axel Sieber s262006}
\date{November 2026}

\begin{document}

\maketitle

\begin{abstract}
% 5-8 sentences: problem, data, method, main numbers, recommendation
\end{abstract}

\tableofcontents

% ======================================================================
\section{Introduction}
% Structure agreed with Tonboe, meeting 2026-11-24:
% Intro -> Theory -> Data -> Methodology -> Results -> Discussion -> Conclusion

\subsection{Passive microwave radiometry of sea ice}
% why microwaves, why sea ice concentration (SIC), long climate record

\subsection{Framework}
% ESMR SIC record (Kolbe et al. 2024), ESA CCI / OSI SAF, RTM noise correction (Tonboe et al.)

\subsection{Where this project fits}
% earliest satellite SIC record (1972-77), one channel only -> noise from the atmosphere and the surface

\subsection{Objectives}
% Task 1: open-water SIC noise, uncorrected vs. corrected
% Task 2: which ERA5 parameters matter (sensitivity)
% Task 3: dependence on incidence angle
% Task 4: recommendation for an ESMR correction strategy

% ======================================================================
\section{Theory}
% Tonboe: "Siri" chapter, assumed to mean theory (Danish "teori")

\subsection{Nimbus 5 ESMR}
% 19.35 GHz, H-pol, cross-track scan, 78 positions, incidence angle 0-64 deg, footprint

\subsection{Sea ice concentration from one channel}
% tie points, SIC = (TB - Tw)/(Ti - Tw)

\subsection{Radiative transfer model and noise correction}
% emission model, V, W, L, Ts, Ti; dTB; iterative correction (Kolbe et al. 2024)

\subsection{Sensitivity of the brightness temperature to the ERA5 inputs}
% RTM sweep (B1), SIC change per 1 sigma; WITHOUT the L study (-> Results)
% wind and the other inputs scale linearly with SIC (water effects largest at c = 0)

% ======================================================================
\section{Data} %Mingming

\subsection{ESMR swath data, February 1973}

\subsection{ERA5 reanalysis}
% reanalysis = summary of all observations that existed in 1973 -> argument and limit

\subsection{Preprocessing and quality filters}
% Kolbe filters Eq. 1-4; outer scan positions 1-4 and 75-78 are KEPT (Tonboe 2026-11-24)

% ======================================================================
\section{Methodology} %Mingming

\subsection{Tie-point selection}
% 5x5 box, pixel lsm = 0, <= 10 % land in the box; water sst > 272 K (near freezing point), |lat| > 55 deg

\subsection{Uncorrected and corrected SIC}
% iterative RTM correction, c clipped to 0..1 inside the RTM only, no 15 % limit; SIC shown unclipped

\subsection{Uncertainty}
% algorithm + resampling uncertainty (Kolbe et al. 2024, Eq. 17-19), 30 km grid

\subsection{Evaluation}
% std (and mean) of open-water SIC, exclusion test, angle bins, case-study selection

% ======================================================================
\section{Results}

\subsection{Mean and STD for the open water SIC for corrected and uncorrected Tbs}
% Task 1

\subsection{Most important parameters for the correction} %sensitivity analysis
% Task 2: exclusion test, data vs. RTM, L study (cluster-mean L vs. per-pixel L)

\subsection{Incidence angle}
% Task 3, including the outer scan positions

\subsection{Uncertainty}
% maps; higher over multiyear ice

\subsection{Case studies}
\subsubsection{A low-pressure system, 12 February 1973}
\subsubsection{Low-pressure systems over the whole month}
\subsubsection{Comparison with the NASA atlases (SP-459, SP-489)}

\subsection{Recommendation for an ESMR correction strategy}
% Task 4

% ======================================================================
\section{Discussion}

\subsection{Geographical difference for the task 1 correction}
% where the correction performs better

\subsection{Geographical difference for the task 2 parameters}
% discuss what parameters are important for the specific region/ hemispheres and due to what reasons
% discuss interannual changes and seasonal changes related to the 1973 Feb (spatial and temporal changes about task 2)

\subsection{Ice edge and multiyear ice}

\subsection{Limitations}
% ERA5 in 1973, cloud liquid L not corrected, one month only

\subsection{Deviations from Kolbe et al. (2024)}
% 272 K water tie point, outer scan positions kept, no 15 % limit

% ======================================================================
\section{Conclusion}
% return to the objectives (Task 1-4), one paragraph each

% ======================================================================
\section*{References}

\appendix
\section{Source code}

\end{document}
